(1) 一束方向与轴线夹角 \(\theta\) 的准单色平面波在相距 \(d\) 的两个狭缝处有光程差约为 \(d\theta\)，因此相位差 \(\Delta\phi(\theta)=kd\theta\)，其中 \(k=2\pi/\lambda\)。不同方向微小点源相互不相干，计算总干涉项必须先对每个方向的复相位因子按强度平均。
(2) 均匀角亮度下复互相干度 \(\gamma_{12}=\frac1{2\theta_0}\int_{-\theta_0}^{\theta_0}e^{ikd\theta}\,d\theta=\sin(kd\theta_0)/(kd\theta_0)\)。记 \(\operatorname{sinc}(x)=\sin x/x\)，则 \(\gamma_{12}=\operatorname{sinc}(kd\theta_0)\)，在 \(x=0\) 处按连续极限取一；该函数来自扩展源角度亮度的傅里叶变换。
(3) 双缝的干涉最大、最小强度一般为 \(I_{\max,\min}=I_1+I_2\pm2\sqrt{I_1I_2}|\gamma_{12}|\)。若 \(I_1=I_2\)，可见度 \(\mathcal V=(I_{\max}-I_{\min})/(I_{\max}+I_{\min})=|\gamma_{12}|\)。第一个零点满足 \(kd\theta_0=\pi\)，从而 \(d_0=\lambda/(2\theta_0)\)，描述在给定源角宽时缝距太大便失去一阶干涉对比。
(4) 当 \(\theta_0\to0\)、\(d\) 固定时 \(\operatorname{sinc}(kd\theta_0)\to1\)，接近单一方向点源的完全空间相干性。对于宽角源、固定大 \(d\)，振荡相位因相互抵消使 \(|\gamma_{12}|\) 整体下降，条纹可能消失，但每个狭缝的平均强度 \(I_1,I_2\) 仍非零；消失的是相干交叉项而不是能量。
